\frfilename{mt266.tex}
\versiondate{28.1.09}
\copyrightdate{2004}

\def\chaptername{Perubahan variabel dalam integral}
\def\sectionname{Ketaksamaan Brunn-Minkowski}

\newsection{*266\dvAnew{2004}}

Sekarang kita telah memiliki sebagian besar unsur penting untuk membuktikan
ketaksamaan Brunn-Minkowski (266C) dalam bentuk yang kuat. Saat ini saya tidak
berharap akan menggunakannya dalam risalah ini, tetapi ketaksamaan tersebut
merupakan salah satu hasil dasar teori ukuran geometris dan, dari posisi kita
sekarang, tidak sulit; karena itu saya menyertakannya di sini. Hasil-hasil
pendahuluan mengenai rataan aritmetika dan geometrik (266A) serta penutupan
esensial (266B) sangat penting pula untuk alasan-alasan lain.

\leader{266A}{Rataan aritmetika dan geometrik \dvrocolon{}}\cmmnt{ Kita
akan memerlukan hasil baku berikut.

\medskip

\noindent}{\bf Proposisi} Jika
$u_0,\ldots,u_n,p_0,\ldots,p_n\in\coint{0,\infty}$ dan
$\sum_{i=0}^np_i=1$, maka $\prod_{i=0}^nu_i^{p_i}\le\sum_{i=0}^np_iu_i$.

\proof{ Lakukan induksi pada $n$. Untuk $n=0$, $p_0=1$, hasilnya
segera. Jika $n=1$, maka apabila $u_1=0$ hasilnya segera (bahkan jika,
sebagaimana menjadi kebiasaan dalam buku ini, kita menafsirkan $0^0$
sebagai $1$). Jika tidak, tetapkan $t=\Bover{u_0}{u_1}$; maka

\Centerline{$t^{p_0}\le p_0t+1-p_0=p_0t+p_1$}

\noindent(seperti pada bagian (a) bukti 244E), sehingga

\Centerline{$u_0^{p_0}u_1^{p_1}=t^{p_0}u_1\le p_0tu_1+p_1u_1
=p_0u_0+p_1u_1$.}

\noindent Untuk langkah induksi menuju $n\ge 2$, jika
$p_0=\ldots=p_{n-1}=0$ hasilnya segera. Jika tidak, tetapkan
$q=p_0+\ldots+p_{n-1}=1-p_n$; maka

$$\eqalignno{\prod_{i=0}^nu_i^{p_i}
&=(\prod_{i=0}^{n-1}u_i^{p_i/q})^qu_n^{p_n}
\le(\sum_{i=0}^{n-1}\Bover{p_i}qu_i)^qu_n^{p_n}\cr
\displaycause{berdasarkan hipotesis induksi}
&\le q(\sum_{i=0}^{n-1}\Bover{p_i}qu_i)+p_nu_n\cr
\displaycause{berdasarkan kasus dua suku yang baru saja diperiksa}
&=\sum_{i=0}^np_iu_i,\cr}$$

\noindent dan induksi berlanjut.
}%end of proof of 266A

\leader{266B}{Proposisi} Untuk setiap himpunan $D\subseteq\BbbR^r$,
tetapkan

\Centerline{$\clstar D=\{x:\limsup_{\delta\downarrow 0}
\Bover{\mu^*(D\cap B(x,\delta))}{\mu B(x,\delta)}>0\}$,}

\noindent dengan $\mu$ ukuran Lebesgue pada $\BbbR^r$.

(a) $D\setminus\clstar D$ terabaikan.

(b) $\clstar D\subseteq\overline{D}$.

(c) $\clstar D$ adalah himpunan Borel.

(d) $\mu(\clstar D)=\mu^*D$.

(e) Jika $C\subseteq\Bbb R^r$, maka
$\overline{C}+\clstar D\subseteq\clstar(C+D)$, dengan $C+D$ menyatakan
$\{x+y:x\in C$, $y\in D\}$.

\proof{{\bf (a)} 261Da.

\medskip

{\bf (b)} Jika $x\in\BbbR^r\setminus\overline{D}$, maka
$D\cap B(x,\delta)=\emptyset$ untuk semua $\delta$ yang cukup kecil.

\medskip

{\bf (c)} Intinya hanyalah bahwa
$(x,\delta)\mapsto\mu^*(D\cap B(x,\delta))$ kontinu. \Prf\ Untuk
setiap $x$, $y\in\BbbR^r$ dan $\delta$, $\eta\ge 0$, kita mempunyai

$$\eqalignno{|\mu^*(D\cap B(y,\eta))-\mu^*(D\cap B(x,\delta))|
&\le\mu(B(y,\eta)\symmdiff B(x,\delta))\cr
&=2\mu(B(x,\delta)\cup B(y,\eta))-\mu B(x,\delta)-\mu B(y,\eta)\cr
&\le\beta_r\bigl(2(\max(\delta,\eta)+\|x-y\|)^r-\delta^r-\eta^r\bigr)\cr
\displaycause{dengan $\beta_r=\mu B(\tbf{0},1)$}
&\to 0\cr}$$

\noindent ketika $(y,\eta)\to(x,\delta)$.\ \QeD\ Jadi

\Centerline{$x\mapsto
\limsup_{\delta\downarrow 0}
  \Bover{\mu^*(D\cap B(x,\delta))}{\mu B(x,\delta)}
=\inf_{\alpha\in\Bbb Q,\alpha>0}\sup_{\beta\in\Bbb Q,0<\beta\le\alpha}
  \Bover1{\beta_r\beta^r}\mu^*(D\cap B(x,\beta))$}

\noindent terukur Borel, dan

\Centerline{$\clstar D=\{x:\limsup_{\delta\downarrow 0}
  \Bover{\mu^*(D\cap B(x,\delta))}{\mu B(x,\delta)}>0\}$}

\noindent adalah himpunan Borel.

\medskip

{\bf (d)} Berdasarkan (c), $\mu(\clstar D)$ terdefinisi; berdasarkan
(a), $\mu(\clstar D)\ge\mu^*D$. Di sisi lain, ambil selubung terukur
$E$ dari $D$ (132Ee); maka 261Db memberi tahu kita bahwa

\Centerline{$\limsup_{\delta\downarrow 0}
  \Bover{\mu^*(D\cap B(x,\delta))}{\mu B(x,\delta)}
\le\limsup_{\delta\downarrow 0}
  \Bover{\mu(E\cap B(x,\delta))}{\mu B(x,\delta)}
=0$}

\noindent untuk hampir setiap $x\in\BbbR^r\setminus E$, sehingga
$\clstar D\setminus E$ terabaikan dan

\Centerline{$\mu(\clstar D)\le\mu E=\mu^*D$.}

\medskip

{\bf (e)} Jika $x\in\overline{C}$ dan $y\in\clstar D$, tetapkan

\Centerline{$\gamma=\Bover13\limsup_{\delta\downarrow 0}
  \Bover{\mu^*(D\cap B(y,\delta))}{\mu B(y,\delta)}>0$.}

\noindent Untuk setiap $\eta>0$, terdapat $\delta\in\ocint{0,\eta}$
sedemikian sehingga
$\mu^*(D\cap B(y,\delta))\ge 2\gamma\mu B(x,\delta)$. Ambil
$\delta_1\in\coint{0,\delta}$ sedemikian sehingga
$\delta^r-\delta_1^r\le\gamma\delta^r$. Maka terdapat $x'\in C$
sedemikian sehingga $\|x-x'\|\le\delta-\delta_1$. Dalam hal ini,

$$\eqalign{\mu^*((C+D)\cap B(x+y,\delta))
&\ge\mu^*((x'+D)\cap B(x'+y,\delta_1))
=\mu^*(D\cap B(y,\delta_1))\cr
&\ge\mu^*(D\cap B(y,\delta))-\mu B(y,\delta)+\mu B(y,\delta_1)\cr
&\ge 2\beta_r\gamma\delta^r-\beta_r\delta^r+\beta_r\delta_1^r
\ge\beta_r\gamma\delta^r.\cr}$$

\noindent Karena $\eta$ sebarang,

\Centerline{$\limsup_{\delta\downarrow 0}
  \Bover{\mu^*((C+D)\cap B(x+y,\delta))}{\mu B(y,\delta)}
\ge\gamma$}

\noindent dan $x+y\in\clstar(C+D)$; karena $x$ dan $y$ sebarang,
$\overline{C}+\clstar D\subseteq\clstar(C+D)$.
}%end of proof of 266B

\cmmnt{\medskip

\noindent{\bf Catatan} Dalam konteks ini, $\clstar D$ disebut
{\bf penutupan esensial} dari $D$.
}

\leader{266C}{Teorema} Misalkan $A$, $B\subseteq\BbbR^r$ adalah
himpunan tak kosong, dengan $r\ge 1$ bilangan bulat. Jika $\mu$ adalah
ukuran Lebesgue pada $\BbbR^r$, dan $A+B=\{x+y:x\in A$, $y\in B\}$,
maka $\mu^*(A+B)^{1/r}\ge(\mu^*A)^{1/r}+(\mu^*B)^{1/r}$.

\proof{{\bf (a)} Pertama-tama tinjau kasus ketika
$A=\coint{a,a'}$ dan $B=\coint{b,b'}$ adalah interval setengah terbuka.
Dalam hal ini $A+B=\coint{a+b,a'+b'}$; dengan menulis
$a=(\alpha_1,\ldots,\alpha_r)$, dan seterusnya, seperti dalam \S115,
tetapkan

\Centerline{$u_i
=\Bover{\alpha'_i-\alpha_i}{\alpha'_i+\beta'_i-\alpha_i-\beta_i}$,
\quad$v_i
=\Bover{\beta'_i-\beta_i}{\alpha'_i+\beta'_i-\alpha_i-\beta_i}$}

\noindent untuk setiap $i$. Maka kita mempunyai

$$\eqalignno{(\mu A)^{1/r}+(\mu B)^{1/r}
&=\prod_{i=1}^r(\alpha'_i-\alpha_i)^{1/r}
  +\prod_{i=1}^r(\beta'_i-\beta_i)^{1/r}\cr
&=\mu(A+B)^{1/r}(\prod_{i=1}^ru_i^{1/r}+\prod_{i=1}^rv_i^{1/r})\cr
&\le\mu(A+B)^{1/r}(\Bover1r\sum_{i=1}^ru_i+\Bover1r\sum_{i=1}^rv_i)\cr
\displaycause{266A}
&=\mu(A+B)^{1/r}.\cr}$$

\medskip

{\bf (b)} Sekarang saya buktikan dengan induksi pada $m+n$ bahwa jika
$A=\bigcup_{j=0}^mA_j$ dan $B=\bigcup_{j=0}^nB_j$, dengan
$\langle A_j\rangle_{j\le m}$ dan $\langle B_j\rangle_{j\le n}$
keduanya keluarga saling lepas dari interval setengah terbuka tak kosong,
maka $\mu(A+B)^{1/r}\ge(\mu A)^{1/r}+(\mu B)^{1/r}$. \Prf\ Induksi
dimulai dengan kasus $m=n=0$, yang telah ditangani dalam (a). Untuk
langkah induksi menuju $m+n=l\ge 1$, salah satu dari $m$, $n$ tidak nol;
argumennya sama dalam kedua kasus; misalkan yang pertama.
Karena $A_0\cap A_1=\emptyset$, harus ada suatu $j\le r$ dan
$\alpha\in\Bbb R$ sedemikian sehingga $A_0$ dan $A_1$ dipisahkan oleh
hiperbidang $\{x:\xi_j=\alpha\}$. Tetapkan
$A'=\{x:x\in A$, $\xi_j<\alpha\}$ dan
$A''=\{x:x\in A$, $\xi_j\ge\alpha\}$; maka $A'$ dan $A''$ keduanya tak
kosong dan dapat dinyatakan sebagai gabungan paling banyak $m$ interval
setengah terbuka yang saling lepas. Tetapkan
$\gamma=\Bover{\mu A'}{\mu A}\in\ooint{0,1}$. Fungsi
$\beta\mapsto\mu\{x:x\in B$, $\xi_j<\beta\}$ kontinu, sehingga terdapat
$\beta\in\Bbb R$ sedemikian sehingga $\mu B'=\gamma\mu B$, dengan
$B'=\{x:x\in B$, $\xi_j<\beta\}$; tetapkan $B''=B\setminus B'$. Maka
$B'$ dan $B''$ dapat dinyatakan sebagai gabungan paling banyak $n+1$
interval setengah terbuka. Berdasarkan hipotesis induksi,

\Centerline{$\mu(A'+B')^{1/r}\ge(\mu A')^{1/r}+(\mu B')^{1/r}$,
\quad$\mu(A''+B'')^{1/r}\ge(\mu A'')^{1/r}+(\mu B'')^{1/r}$.}

\noindent Sekarang $A'+B'\subseteq\{x:\xi_j<\alpha+\beta\}$,
sedangkan $A''+B''\subseteq\{x:\xi_j\ge\alpha+\beta\}$. Jadi

$$\eqalign{\mu(A+B)
&\ge\mu(A'+B')+\mu(A''+B'')\cr
&\ge\bigl((\mu A')^{1/r}+(\mu B')^{1/r}\bigr)^r
  +\bigl((\mu A'')^{1/r}+(\mu B'')^{1/r}\bigr)^r\cr
&=\bigl((\gamma\mu A)^{1/r}+(\gamma\mu B)^{1/r}\bigr)^r
  +\bigl(((1-\gamma)\mu A)^{1/r}+((1-\gamma)\mu B)^{1/r}\bigr)^r\cr
&=(\bigl(\mu A)^{1/r}+(\mu B)^{1/r}\bigr)^r.\cr}$$

\noindent Dengan mengambil akar ke-$r$,
$\mu(A+B)^{1/r}\ge(\mu A)^{1/r}+(\mu B)^{1/r}$ dan induksi
berlanjut.\ \Qed

\medskip

{\bf (c)} Sekarang misalkan $A$ dan $B$ adalah himpunan bagian kompak
tak kosong dari $\BbbR^r$. Maka
$\mu(A+B)^{1/r}\ge(\mu A)^{1/r}+(\mu B)^{1/r}$.
\Prf\ $A+B$ kompak (karena $A\times B\subseteq\BbbR^r\times\BbbR^r$
kompak, sebab tertutup dan terbatas, serta penjumlahan kontinu, sehingga
kita dapat menggunakan 2A2Eb). Ambil $\epsilon>0$. Ambil himpunan terbuka
$G\supseteq A+B$ sedemikian sehingga $\mu G\le\mu(A+B)+\epsilon$
(134Fa); maka terdapat $\delta>0$ sedemikian sehingga
$B(x,2\delta)\subseteq G$ untuk setiap $x\in A+B$ (2A2Ed). Ambil
$n\in\Bbb N$ sedemikian sehingga $2^{-n}\sqrt{r}\le\delta$, dan misalkan
$A_1$ adalah gabungan semua interval setengah terbuka berbentuk
$\coint{2^{-n}z,2^{-n}z+2^{-n}e}$ yang bertemu $A$, dengan
$z\in\BbbZ^r$ dan $e=(1,1,\ldots,1)$. Maka $A_1$ adalah gabungan
berhingga yang saling lepas dari interval-interval setengah terbuka,
$A\subseteq A_1$, dan setiap titik dalam $A_1$ berjarak paling jauh $\delta$
dari suatu titik dalam $A$. Dengan cara serupa, kita dapat memperoleh himpunan
$B_1$, suatu gabungan berhingga yang saling lepas dari interval-interval
setengah terbuka, yang mencakup $B$ dan sedemikian sehingga setiap titik
dalam $B_1$ berjarak paling jauh $\delta$ dari suatu titik dalam $B$. Namun ini
berarti bahwa setiap titik dalam $A_1+B_1$ berjarak paling jauh $2\delta$ dari
suatu titik dalam $A+B$, dan termasuk dalam $G$. Oleh karena itu

$$\eqalignno{(\mu(A+B)+\epsilon)^{1/r}
&\ge(\mu G)^{1/r}
\ge\mu(A_1+B_1)^{1/r}
\ge(\mu A_1)^{1/r}+(\mu B_1)^{1/r}\cr
\displaycause{berdasarkan (b)}
&\ge(\mu A)^{1/r}+(\mu B)^{1/r}.\cr}$$

\noindent Karena $\epsilon$ sebarang,
$\mu(A+B)^{1/r}\ge(\mu A)^{1/r}+(\mu B)^{1/r}$.\ \Qed

\medskip

{\bf (d)} Selanjutnya, misalkan $A$, $B\subseteq\BbbR^r$ terukur
Lebesgue. Maka

$$\eqalignno{(\mu A)^{1/r}+(\mu B)^{1/r}
&=\sup\{(\mu K)^{1/r}+(\mu L)^{1/r}:
  K\subseteq A\text{ dan }L\subseteq B\text{ kompak}\}\cr
\displaycause{134Fb}
&\le\sup\{\mu(K+L)^{1/r}:
  K\subseteq A\text{ dan }L\subseteq B\text{ kompak}\}\cr
\displaycause{berdasarkan (c)}
&\le\mu^*(A+B)^{1/r}.\cr}$$

\medskip

{\bf (e)} Untuk langkah kedua terakhir, misalkan
$A$, $B\subseteq\BbbR^r$ mempunyai ukuran luar Lebesgue tak nol. Tinjau
$\clstar A$, $\clstar B$, dan $\clstar(A+B)$ sebagaimana didefinisikan
dalam 266B. Maka $\clstar A$ dan $\clstar B$ tak kosong dan jumlahnya
termuat dalam $\clstar(A+B)$, berdasarkan 266Bb dan 266Be. Jadi kita
mempunyai

$$\eqalignno{(\mu^*A)^{1/r}+(\mu^*B)^{1/r}
&=\mu(\clstar A)^{1/r}+\mu(\clstar B)^{1/r}\cr
\displaycause{266Bd}
&\le\mu^*(\clstar A+\clstar B)^{1/r}\cr
\displaycause{berdasarkan (d) di sini}
&\le\mu(\clstar(A+B))^{1/r}
=\mu^*(A+B)^{1/r}.\cr}$$

\medskip

{\bf (f)} Akhirnya, untuk himpunan tak kosong sebarang $A$,
$B\subseteq\BbbR^r$, perhatikan bahwa jika (misalnya) $A$ dapat
diabaikan, kita dapat mengambil sembarang $x\in A$ dan melihat bahwa

\Centerline{$\mu^*(A+B)^{1/r}\ge\mu^*(x+B)^{1/r}=(\mu^*B)^{1/r}
=(\mu^*A)^{1/r}+(\mu^*B)^{1/r}$,}

\noindent dan hasilnya segera pula jika $B$ terabaikan. Jadi semua
kasus telah tercakup.
}%end of proof of 266C

\exercises{\leader{266X}{Latihan dasar (a)}
Misalkan $D$, $D'$ adalah himpunan bagian dari $\BbbR^r$. Tunjukkan bahwa
(i) $\clstar(D\cup D')=\clstar D\cup\clstar D'$ (ii) $\clstar D=\clstar D'$
jika dan hanya jika $D$ dan $D'$ mempunyai selubung terukur bersama
(iii) $\clstar D\setminus\clstar(\BbbR^r\setminus D')
\subseteq\clstar(D\cap D')$ (iv) $D$ terukur Lebesgue jika dan hanya
jika $\clstar D\cap\clstar(\BbbR^r\setminus D)$ terabaikan menurut
Lebesgue (v) $D\cup\clstar D$ adalah selubung terukur dari $D$
(vi) $\clstar(\clstar D)=\clstar D$.
%266B

\spheader 266Xb Tunjukkan bahwa, untuk himpunan terukur
$E\subseteq\Bbb R$, $\clstar E$ tepat merupakan himpunan bilangan real
yang bukan titik kepadatan dari $\Bbb R\setminus E$.
%266B

\spheader 266Xc Dalam 266C, tunjukkan bahwa jika $A$ dan $B$ adalah
himpunan cembung serupa dengan orientasi yang sama, maka $A+B$ adalah
himpunan cembung yang serupa dengan keduanya dan
$\mu(A+B)^{1/r}=(\mu A)^{1/r}+(\mu B)^{1/r}$.
%266C

%\leader{266Y}{Further exercises (a)}

}%end of exercises

\endnotes{
\Notesheader{266}
Bukti 266C diambil dari {\smc Federer 69}.
%Is there a Hausdorff-measure version?
Ada sedikit generalitas semu dalam bentuk yang diberikan di sini. Jika
himpunan $A$ dan $B$ sedikit saja tak beraturan, maka
$\mu^*(A+B)^{1/r}$ kemungkinan jauh lebih besar daripada
$(\mu^*A)^{1/r}+(\mu^*B)^{1/r}$. Kasus kritis, ketika $A$ dan $B$
merupakan himpunan cembung serupa, jauh lebih mudah (266Xc). Karena itu
teorema ini paling berguna ketika $A$ dan $B$ adalah himpunan cembung
yang tidak serupa dan kita memperoleh taksiran tak sepele yang mungkin
sulit dibuktikan dengan cara lain. Untuk kasus ini kita tidak memerlukan
266B. Teorema 266C merupakan contoh yang mencerahkan tentang cara dimensi
$r$ masuk ke dalam rumus-rumus ketika kita mencari hasil yang berlaku
bagi ruang Euklides umum. Akan ada jauh lebih banyak contoh ketika saya
kembali ke teori ukuran geometris dalam Bab 47 Jilid 4.
}%end of notes

\discrpage

